\documentclass[12pt, a4paper]{article}

% ── Packages ──────────────────────────────────────────────────────────────────
\usepackage[margin=2.5cm]{geometry}
\usepackage{amsmath, amssymb, amsthm}
\usepackage{booktabs}
\usepackage{array}
\usepackage{multirow}
\usepackage{tabularx}
\usepackage{longtable}
\usepackage[table]{xcolor}
\usepackage{graphicx}
\usepackage{hyperref}
\usepackage{enumitem}
\usepackage{parskip}
\usepackage{titlesec}
\usepackage{fancyhdr}
\usepackage{caption}
\usepackage{tcolorbox}
\usepackage{mdframed}
\usepackage{microtype}
\usepackage{url}
\usepackage[T1]{fontenc}
\usepackage{lmodern}

% ── Colours ───────────────────────────────────────────────────────────────────
\definecolor{gapblue}{RGB}{13, 71, 161}
\definecolor{gapgold}{RGB}{230, 162, 0}
\definecolor{gapgreen}{RGB}{27, 94, 32}
\definecolor{gapred}{RGB}{183, 28, 28}
\definecolor{lightgray}{RGB}{245, 245, 245}
\definecolor{headerblue}{RGB}{227, 242, 253}
\definecolor{ideabg}{RGB}{232, 245, 233}
\definecolor{warnbg}{RGB}{255, 243, 224}
\definecolor{critbg}{RGB}{255, 235, 238}

% ── Section formatting ────────────────────────────────────────────────────────
\titleformat{\section}{\large\bfseries\color{gapblue}}{{\thesection}}{1em}{}[\titlerule]
\titleformat{\subsection}{\normalsize\bfseries\color{gapblue!80!black}}{\thesubsection}{1em}{}
\titleformat{\subsubsection}{\normalsize\bfseries\itshape}{\thesubsubsection}{1em}{}

% ── tcolorbox styles ──────────────────────────────────────────────────────────
\tcbuselibrary{skins, breakable}

\newtcolorbox{ideabox}[2][]{
  colback=ideabg, colframe=gapgreen!70!black,
  fonttitle=\bfseries\color{gapgreen!30!black},
  title={#2}, #1, breakable,
  left=6pt, right=6pt, top=4pt, bottom=4pt
}

\newtcolorbox{gapbox}[2][]{
  colback=headerblue, colframe=gapblue!70!black,
  fonttitle=\bfseries\color{gapblue!30!black},
  title={#2}, #1, breakable,
  left=6pt, right=6pt, top=4pt, bottom=4pt
}

\newtcolorbox{warnbox}[1][]{
  colback=warnbg, colframe=gapgold!80!black,
  fonttitle=\bfseries, #1, breakable,
  left=6pt, right=6pt, top=4pt, bottom=4pt
}

\newtcolorbox{critbox}[1][]{
  colback=critbg, colframe=gapred!70!black,
  fonttitle=\bfseries\color{gapred!30!black}, #1, breakable,
  left=6pt, right=6pt, top=4pt, bottom=4pt
}

% ── Header / footer ───────────────────────────────────────────────────────────
\pagestyle{fancy}
\fancyhf{}
\fancyhead[L]{\small\textcolor{gapblue}{\textbf{REX-SUB Gap Analysis}}}
\fancyhead[R]{\small\textcolor{gray}{Research Extension Report}}
\fancyfoot[C]{\small\thepage}
\renewcommand{\headrulewidth}{0.4pt}

% ── Hyperref setup ────────────────────────────────────────────────────────────
\hypersetup{
  colorlinks=true, linkcolor=gapblue, citecolor=gapblue, urlcolor=gapblue,
  pdftitle={REX-SUB: Research Gap Analysis and Extension Directions},
  pdfauthor={Auto-Research Skills Analysis}
}

% ── Custom commands ───────────────────────────────────────────────────────────
\newcommand{\gap}[1]{\textcolor{gapred}{\textbf{Gap:} #1}}
\newcommand{\highlight}[1]{\textcolor{gapblue}{\textbf{#1}}}
\newcommand{\novelty}[1]{\textbf{Novelty:} #1/5}
\newcommand{\feasibility}[1]{\textbf{Feasibility:} #1/5}
\newcommand{\mspe}{$\phi_{\text{MSPE}}$}
\newcommand{\IS}{$\phi_{\text{IS}}$}
\newcommand{\Vecchia}{\textsc{Vecchia}}

% ─────────────────────────────────────────────────────────────────────────────
\begin{document}

% ── Title page ────────────────────────────────────────────────────────────────
\begin{titlepage}
  \centering
  \vspace*{2cm}
  {\Huge\bfseries\color{gapblue}
    REX-SUB: Research Gap Analysis\\[0.4em]
    \& Extension Directions\par}
  \vspace{1.5em}
  {\large\itshape
    A systematic analysis of \emph{REX-SUB: A scalable subsampling strategy\\
    for modeling large spatial datasets}\\[0.3em]
    Rios \& Lee, \emph{Computational Statistics and Data Analysis}
    224 (2026) 108404\par}
  \vspace{2em}
  \hrule height 0.8pt
  \vspace{1em}
  {\normalsize
    Produced using the \textbf{AI Auto-Research Skill Suite}\\
    Skills applied: \texttt{/literature-review} \quad
                    \texttt{/peer-review} \quad
                    \texttt{/idea-generation}\par}
  \vspace{0.5em}
  {\small\color{gray}Based on: Kong et al.\ (2026).
    \emph{AI for Auto-Research: Roadmap \& User Guide.}
    arXiv:2605.18661v1\par}
  \vspace{1em}
  \hrule height 0.8pt
  \vspace{2em}
  {\small Generated: June 2026}
  \vfill
  \begin{tcolorbox}[colback=headerblue, colframe=gapblue, width=0.9\textwidth]
    \centering\small
    This document applies a structured AI-assisted research lifecycle framework
    to identify gaps in the REX-SUB paper and generate specific,
    scored, and feasibility-assessed research extension hypotheses.
  \end{tcolorbox}
\end{titlepage}

% ── Table of Contents ─────────────────────────────────────────────────────────
\tableofcontents
\newpage

% ─────────────────────────────────────────────────────────────────────────────
\section{Paper Overview and Core Message}
% ─────────────────────────────────────────────────────────────────────────────

\subsection{Bibliographic Information}

\begin{tabular}{@{}ll@{}}
  \toprule
  \textbf{Title}   & REX-SUB: A scalable subsampling strategy for modeling large spatial datasets \\
  \textbf{Authors} & Nicholas Rios, Ben Seiyon Lee \\
  \textbf{Journal} & Computational Statistics and Data Analysis, 224 (2026) 108404 \\
  \textbf{DOI}     & \texttt{10.1016/j.csda.2026.108404} \\
  \textbf{Keywords}& Design of experiments, Spatial statistics, Optimal subsampling \\
  \bottomrule
\end{tabular}

\subsection{Core Message Card}

\begin{gapbox}{Paper Core Message}
\begin{description}[leftmargin=3em, labelwidth=2.5em]
  \item[\textbf{Problem}]
    Standard Gaussian process (GP) fitting is $\mathcal{O}(N^3)$ in the number of
    spatial locations $N$, making it computationally infeasible for massive spatial
    datasets ($N > 10^4$, e.g.\ satellite data with $N = 2{,}748{,}620$).

  \item[\textbf{Approach}]
    REX-SUB: a \emph{Randomized Exchange algorithm for Spatial SUBsampling} that
    iteratively selects a small subsample of size $n \ll N$ by exchanging locations
    to minimise prediction-based optimality criteria
    (\mspe{} or \IS{}), using the \Vecchia{} approximation for fast GP
    re-fitting at each iteration.

  \item[\textbf{Key result}]
    REX-SUB achieves lower MSPE and interval scores than random subsampling, Latin
    Hypercube Sampling (LHS), and sequential IMSPE-based design (hetGP) across all
    16 simulation scenarios and on the MODIS precipitable water dataset
    ($N = 2{,}748{,}620$) for $n \in \{25, 36\}$.

  \item[\textbf{Who benefits}]
    Environmental statisticians, geostatisticians, and data scientists fitting GP
    models to remotely sensed or sensor-network data at million-location scales.

  \item[\textbf{Important caveat}]
    Assumes a \emph{stationary, isotropic} Matérn covariance; tested only for
    $n \in \{25, 36\}$; no theoretical convergence guarantees; single time-point
    spatial data only.
\end{description}
\end{gapbox}

% ─────────────────────────────────────────────────────────────────────────────
\section{Literature Synthesis (\texttt{/literature-review})}
% ─────────────────────────────────────────────────────────────────────────────

We organise the existing literature into six thematic clusters and assess how
REX-SUB is positioned relative to each.

% ── Cluster A ─────────────────────────────────────────────────────────────────
\subsection{Cluster A: Scalable GP Approximations (Modify the Model)}

\textbf{Representative works:} Vecchia (1988a), Datta et al.\ (2016) [NNGP],
Katzfuss (2017) [MRA], Furrer et al.\ (2006) [tapering], Cressie \& Johannesson
(2008) [fixed-rank kriging], Banerjee et al.\ (2008) [predictive processes],
Katzfuss et al.\ (2020).

\textbf{Shared assumption:} retain all $N$ observations but approximate the
covariance or likelihood structure to achieve sub-cubic computational cost.  These
methods trade approximation fidelity for scalability by modifying the model, not the
data.

\textbf{Relationship to REX-SUB:} orthogonal.  REX-SUB uses the \Vecchia{}
approximation as an \emph{ingredient} within each exchange step rather than treating
it as a competitor.  The two families are complementary: scalable GPs model the full
dataset with approximate likelihoods; REX-SUB models a small, optimally chosen
dataset with an exact (Vecchia-approximated) likelihood.

\begin{warnbox}
  \textbf{Remaining gap:} No work jointly optimises (a) which subsample to select
  \emph{and} (b) which GP approximation to use.  Could REX-SUB choose subsamples
  that improve the conditioning of \Vecchia{} conditioning sets, reducing
  approximation error simultaneously with prediction error?
\end{warnbox}

% ── Cluster B ─────────────────────────────────────────────────────────────────
\subsection{Cluster B: Spatial Subsampling (Reduce the Data)}

\textbf{Representative works:} Hayashi et al.\ (2020) [random subsampling theory],
McKay et al.\ (1979) \& Zhao et al.\ (2018) [LHS], Heaton et al.\ (2019)
[large-spatial-data competition], Binois \& Gramacy (2021) [hetGP/IMSPE],
Saha \& Bradley (2024) [Bayesian data subset model], Lee \& Park (2023).

\textbf{Shared motivation:} reduce the data size prior to GP fitting to avoid
costly matrix operations.  Most methods rely on geometric space-filling heuristics
(LHS, domain partitioning) or block-independence assumptions rather than
prediction-based optimality.

\textbf{Key tension:} hetGP (Binois \& Gramacy, 2021) uses a prediction-based
criterion (IMSPE) but requires known covariance parameters as input, limiting its
use to iterative sequential settings.  REX-SUB avoids this by simultaneously
estimating parameters and optimising the subsample.  Saha \& Bradley (2024) embed
subsampling in a Bayesian hierarchical framework but provide no optimality criterion
for prediction accuracy.

\begin{warnbox}
  \textbf{Remaining gap:} No method provides principled guidance on the choice of
  subsample size $n$.  All methods, including REX-SUB, require $n$ to be
  pre-specified by the user.
\end{warnbox}

% ── Cluster C ─────────────────────────────────────────────────────────────────
\subsection{Cluster C: Optimal Subsampling for Non-Spatial Models}

\textbf{Representative works:} Wang et al.\ (2019) [IBOSS -- D-optimality for
linear regression], Cheng et al.\ (2020) [logistic IBOSS], Singh \& Stufken (2023)
[LASSO], Liu et al.\ (2023) [cluster-wise regression], Shi \& Tang (2021)
[maximin-robust subsampling].

\textbf{Shared assumption:} independent and identically distributed (IID) or
exchangeable observations.  The D-optimality criterion (Fisher information
maximisation) is used because it leads to closed-form solutions and asymptotic
efficiency guarantees.

\textbf{REX-SUB's contribution:} first extension of prediction-optimised subsampling
into the spatially correlated GP setting, where IID assumptions fail and D-optimality
cannot be directly applied because the GP likelihood couples all observations through
the covariance matrix.

\begin{warnbox}
  \textbf{Remaining gap:} IBOSS provides asymptotic efficiency theory and
  consistency results.  REX-SUB provides neither.  The bridge between optimal
  experimental design theory and spatial GP subsampling is entirely absent.
\end{warnbox}

% ── Cluster D ─────────────────────────────────────────────────────────────────
\subsection{Cluster D: Non-Stationary Spatial Models}

\textbf{Representative works:} Fuentes (2001), Kim et al.\ (2005) [locally
stationary], Risser \& Calder (2015) [convoSPAT], Hefley et al.\ (2017) [basis
functions], Katzfuss (2017) [MRA], Lee \& Park (2023) [scalable partitioned model].

\textbf{Shared motivation:} many real spatial processes exhibit covariance parameters
that vary across the domain (e.g.\ ocean temperature varies with latitude; urban
versus rural aerosol optical depth shows different spatial scales).

\textbf{REX-SUB's assumption:} a single stationary, isotropic Matérn covariance
$K(\mathbf{s}, \mathbf{s}'; \boldsymbol{\Theta})$ with global parameters
$\boldsymbol{\Theta} = (\nu, \phi, \sigma^2, \tau^2)'$.  This is the most
restrictive assumption in the paper.

\begin{critbox}[title=\textbf{Critical Gap (Most Restrictive Assumption)}]
  Non-stationarity fundamentally changes which locations are informative.  In a
  non-stationary field, high-variation regions demand denser subsampling; smooth
  regions need fewer points.  REX-SUB's global GP criterion cannot detect or respond
  to local spatial structure.  The MODIS precipitable water field shown in Figure~2
  of the paper exhibits strong spatial heterogeneity (smooth Atlantic Ocean vs.\
  complex land--ocean boundary), yet the method assumes a single spatial range
  parameter.
\end{critbox}

% ── Cluster E ─────────────────────────────────────────────────────────────────
\subsection{Cluster E: Spatio-Temporal Models}

\textbf{Representative works:} Gneiting (2002), Cressie \& Huang (1999)
[nonseparable covariance functions], Cressie \& Wikle (2011) [textbook],
Zammit-Mangion \& Cressie (2021) [FRK-ST].

\textbf{Shared motivation:} spatial data collected repeatedly over time requires
models that capture both spatial and temporal dependence jointly.

\textbf{REX-SUB's scope:} single spatial snapshot only.  The MODIS dataset used is
from a single day (June 22, 2023), despite the satellite collecting data daily.

\begin{warnbox}
  \textbf{Remaining gap:} No optimal subsampling strategy exists for
  spatio-temporal GPs.  Selecting which (location, time) pairs to include in a
  subsample for spatio-temporal GP fitting is an entirely open problem with
  significantly higher dimensionality.
\end{warnbox}

% ── Cluster F ─────────────────────────────────────────────────────────────────
\subsection{Cluster F: Non-Gaussian Spatial Responses}

\textbf{Representative works:} Zilber \& Katzfuss (2021) [Vecchia-Laplace for
non-Gaussian spatial data], Wood (2004) [GAMs], Lee \& Park (2023) [non-Gaussian
hierarchical models].

\textbf{REX-SUB's scope:} Gaussian responses only.  The precipitable water data are
log-transformed to approximate Gaussianity.  Many spatial datasets (species counts,
disease incidence, habitat presence/absence) cannot be usefully Gaussianised.

\begin{warnbox}
  \textbf{Remaining gap:} The MSPE and Interval Score criteria are derived for
  Gaussian responses.  For non-Gaussian likelihoods, proper scoring rules (log-score,
  CRPS) must replace MSPE.  No prediction-optimal subsampling exists for latent GP
  models.
\end{warnbox}

% ── Master gap table ──────────────────────────────────────────────────────────
\subsection{Consolidated Gap Summary}

\begin{table}[ht]
\centering
\caption{Identified research gaps in REX-SUB and the broader scalable spatial
  statistics literature.  Severity: \textcolor{gapred}{\textbf{H}} = High,
  \textcolor{gapgold}{\textbf{M}} = Medium,
  \textcolor{gapgreen}{\textbf{L}} = Low.}
\label{tab:gaps}
\small
\setlength{\tabcolsep}{5pt}
\begin{tabular}{@{}clp{4.2cm}p{4.2cm}c@{}}
\toprule
\textbf{\#} & \textbf{Gap} & \textbf{Prior work addressing it} &
  \textbf{What is missing} & \textbf{Sev.} \\
\midrule
G1 & Non-stationary covariance
   & convoSPAT, MRA, SPDE, Lee \& Park (2023)
   & Subsampling strategy for non-stationary GPs
   & \textcolor{gapred}{\textbf{H}} \\
\addlinespace[2pt]
G2 & Spatio-temporal extension
   & FRK-ST, Gneiting (2002)
   & Optimal subsampling for space--time GPs
   & \textcolor{gapred}{\textbf{H}} \\
\addlinespace[2pt]
G3 & Expensive GP re-fit per exchange
   & Rank-1 Cholesky updates (general)
   & Fast Vecchia precision-matrix updates on location swap
   & \textcolor{gapred}{\textbf{H}} \\
\addlinespace[2pt]
G4 & Local optima (greedy exchange)
   & Simulated annealing, threshold accepting
   & Global optimisation for spatial subsampling
   & \textcolor{gapgold}{\textbf{M}} \\
\addlinespace[2pt]
G5 & Optimal subsample size $n$
   & Cross-validation (general)
   & Theory or method for adaptive $n$ in spatial GP setting
   & \textcolor{gapgold}{\textbf{M}} \\
\addlinespace[2pt]
G6 & Non-Gaussian responses
   & Vecchia-Laplace (Zilber \& Katzfuss, 2021)
   & CRPS/log-score optimal subsampling for latent GPs
   & \textcolor{gapgold}{\textbf{M}} \\
\addlinespace[2pt]
G7 & Multivariate/co-kriging
   & Co-kriging, multivariate GP literature
   & Joint optimal subsampling for multiple spatial responses
   & \textcolor{gapgold}{\textbf{M}} \\
\addlinespace[2pt]
G8 & Theoretical guarantees
   & IBOSS theory (Wang et al., 2019)
   & Convergence rates and consistency for REX-SUB
   & \textcolor{gapred}{\textbf{H}} \\
\addlinespace[2pt]
G9 & Bias in $\hat{\boldsymbol{\Theta}}$ from informative subsample
   & Informative sampling literature (general)
   & Bias characterisation for GP parameters under optimal subsampling
   & \textcolor{gapgold}{\textbf{M}} \\
\addlinespace[2pt]
G10 & Online/streaming subsampling
   & Online learning (general)
   & Incremental REX-SUB as new spatial data arrives
   & \textcolor{gapgreen}{\textbf{L}} \\
\bottomrule
\end{tabular}
\end{table}

% ─────────────────────────────────────────────────────────────────────────────
\section{Critical Review (\texttt{/peer-review})}
% ─────────────────────────────────────────────────────────────────────────────

\subsection{Paper Summary}

REX-SUB proposes a randomised exchange algorithm for selecting prediction-optimal
subsamples from large spatial datasets, embedding the \Vecchia{} approximation for
scalable GP inference.  Evidence is provided from 16 simulation scenarios covering
combinations of Matérn smoothness $\nu \in \{0.5, 1.5\}$, effective range
$\rho^* \in \{0.3, 0.6\}$, noise proportion $\tau^2/(\tau^2 + \sigma^2) \in
\{0.01, 0.10\}$, and subsample size $n \in \{25, 36\}$, plus one real MODIS
application ($N = 2{,}748{,}620$).

\subsection{Strengths}

\begin{enumerate}[label=\textbf{S\arabic*.}, leftmargin=2.5em]
  \item \textbf{Compelling real application.}  Finding a near-optimal subsample of
    size $n=25$ from 2.7~million locations in $\approx 7$ minutes versus $>1$~hour
    for full-data Vecchia fitting is a practically meaningful result.

  \item \textbf{Elegant Vecchia integration.}  Embedding the \Vecchia{}
    approximation inside each exchange step is a non-trivial engineering choice that
    avoids modifying the GP model while enabling fast per-step evaluation.

  \item \textbf{Comprehensive simulation design.}  The full factorial over
    $\nu \times \rho^* \times \tau^2/(\tau^2+\sigma^2) \times n$ with 100 replicates
    per cell is systematic and reproducible; standard errors are consistently
    reported.

  \item \textbf{Two complementary criteria.}  Jointly assessing MSPE (point
    prediction quality) and Interval Score (uncertainty quantification quality)
    is more rigorous than using MSPE alone.
\end{enumerate}

\subsection{Weaknesses}

\begin{enumerate}[label=\textbf{W\arabic*.}, leftmargin=2.5em]
  \item \textbf{[Decisive] Missing comparison against full-data scalable GP at
    equivalent computational budget.}  The paper does not include an experiment
    comparing REX-SUB ($n=25$, ${\sim}75$\,s) against a full-data NNGP or \Vecchia{}
    fit truncated to the same 75-second budget.  Without this, the central practical
    claim --- that subsampling outperforms scalable approximations at fixed compute
    --- is unsubstantiated.

  \item \textbf{[Major] Subsample sizes $n \in \{25, 36\}$ are unrealistically
    small.}  For $N = 12{,}500$, $n=25$ is $0.2\%$ of the data; for $N=2{,}748{,}620$,
    it is $0.0009\%$.  No MSPE-vs-$n$ tradeoff curve is provided.  Practitioners
    cannot generalise these results to any other $n$.

  \item \textbf{[Major] Stationarity assumption untested under violation.}  All 16
    simulation settings use a stationary isotropic Matérn process --- the exact model
    REX-SUB assumes.  No experiment evaluates performance when the true process is
    non-stationary, despite strong spatial heterogeneity being visible in the MODIS
    application (Figure~2 of the paper).

  \item \textbf{[Major] No theoretical analysis.}  REX-SUB is a greedy heuristic
    with no convergence guarantees, no consistency results, and no bounds on the
    gap between REX-SUB predictions and oracle (full-data GP) predictions.  IBOSS
    --- the closest analogue for non-spatial models --- provides asymptotic efficiency
    theory \cite{Wang2019}.

  \item \textbf{[Minor] Single application domain.}  One NASA satellite dataset on
    one day and one geographic region is insufficient evidence for the claimed broad
    applicability.  A second application (e.g.\ species count surveys or subsurface
    geology) would strengthen generalisability.

  \item \textbf{[Minor] Parallelisation mentioned but not benchmarked.}
    The text states that $n_\text{cand}$ evaluations can be parallelised, but
    no parallel timing results are reported.  Table~5 CPU times therefore represent
    the worst-case sequential implementation.
\end{enumerate}

\subsection{Questions for Future Research}

\begin{enumerate}[label=\textbf{Q\arabic*.}, leftmargin=2.5em]
  \item At what subsample size $n$ does REX-SUB's advantage over LHS vanish?  Is
    there a crossover point where the exchange computation is not worth the MSPE
    gain?

  \item How does REX-SUB perform when the true process is non-stationary?  Does it
    preferentially select locations from smooth or rough regions?

  \item Can the \Vecchia{} precision matrix $\mathbf{Q} = \mathbf{U}\mathbf{U}^\top$
    be updated via a rank-1 (or rank-$m$) modification when a single location is
    swapped, avoiding a full GP re-fit at each exchange?

  \item Are covariance parameter estimates $\hat{\boldsymbol{\Theta}}$ from the
    informatively selected subsample biased?  Does this bias propagate to prediction
    intervals?

  \item Can REX-SUB be extended to multiple time points of MODIS data?  What
    changes in Algorithm~1 when candidates are (location, time) pairs?
\end{enumerate}

% ─────────────────────────────────────────────────────────────────────────────
\section{Research Extension Ideas (\texttt{/idea-generation})}
% ─────────────────────────────────────────────────────────────────────────────

Six specific, scored, and feasibility-assessed research hypotheses are presented
below, ordered by strategic priority.

% ── Idea 1 ───────────────────────────────────────────────────────────────────
\subsection{Idea 1: Efficient REX-SUB via Vecchia Rank-1 Updates}

\begin{ideabox}{Idea~1 \textbullet{} Algorithmic Efficiency [Gap~G3]}

\begin{description}[leftmargin=5em, labelwidth=4.5em]
  \item[\textbf{Hypothesis}]
    Deriving closed-form rank-1 updates to the \Vecchia{} precision matrix
    $\mathbf{Q}$ when a single subsample location is swapped will reduce REX-SUB's
    dominant computational cost by an order of magnitude, making subsample sizes
    $n \geq 200$ practical.

  \item[\textbf{Motivation}]
    Table~5 of the paper shows that 99\% of REX-SUB's CPU time is spent in Steps~2
    and 5 of Algorithm~1 --- fitting GPs and evaluating the \mspe{} criterion.
    Each exchange modifies exactly \emph{one} location while the remaining $n-1$
    locations and their conditioning sets are unchanged.  The \Vecchia{} precision
    matrix $\mathbf{Q} = \mathbf{U}\mathbf{U}^\top$ (where $\mathbf{U}$ is an
    upper-triangular Cholesky factor) can potentially be updated via a rank-$m$
    modification when location $\mathbf{s}_j$ is replaced by $\mathbf{s}_{j'}$,
    analogous to classical Cholesky rank-1 updates in numerical linear algebra.
    A factor-of-$n$ speedup would make REX-SUB applicable to $n \in \{100, 500\}$
    within the same 75-second budget.

  \item[\textbf{Approach}]
    (i) Derive the algebraic form of $\Delta\mathbf{Q}$ when location
    $\mathbf{s}_j$ is replaced in a maxmin-ordered \Vecchia{} subsample.
    (ii)~Identify which rows/columns of $\mathbf{U}$ change --- only those
    involving $\mathbf{s}_j$ or its downstream dependents in the directed acyclic
    graph $G = (S, N_S)$.  (iii)~Implement in C++ via the \texttt{GpGp}
    package infrastructure and wrap in R.

  \item[\textbf{Key experiment}]
    Benchmark REX-SUB (standard) vs.\ REX-SUB (rank-1 update) on the MODIS dataset
    for $n \in \{25, 100, 500\}$.  Report: speedup ratio, MSPE equivalence (updates
    should be mathematically exact), and the resulting MSPE-vs-$n$ curve.

  \item[\textbf{Success criterion}]
    Rank-1 update reduces per-exchange computation by $\geq 10\times$ while
    producing identical MSPE to the full re-fit implementation across all
    tested settings.

  \item[\novelty{4}]
    Rank-1 Cholesky updates are classical; applying them to \Vecchia{} precision
    matrices in the exchange-subsampling context is new.

  \item[\feasibility{4}]
    Algebraically tractable; requires careful treatment of the \Vecchia{} DAG
    structure when a location is swapped.
\end{description}

\end{ideabox}

% ── Idea 2 ───────────────────────────────────────────────────────────────────
\subsection{Idea 2: Adaptive Subsample Size Selection}

\begin{ideabox}{Idea~2 \textbullet{} Practical Guidance [Gap~G5]}

\begin{description}[leftmargin=5em, labelwidth=4.5em]
  \item[\textbf{Hypothesis}]
    A cross-validation procedure tracking marginal MSPE reduction as $n$ increases
    within REX-SUB will identify a data-adaptive elbow point, eliminating the need
    to pre-specify $n$ while providing the first principled guidance on subsample
    size for spatial GP models.

  \item[\textbf{Motivation}]
    REX-SUB requires $n$ to be pre-specified; the paper tests only $n \in \{25, 36\}$
    for $N \in \{10{,}000;\, 2{,}748{,}620\}$.  In practice, the optimal $n$ depends
    on: the spatial range $\phi$ (more informative observations when $\phi$ is large),
    the noise ratio $\tau^2/(\tau^2+\sigma^2)$ (higher noise requires larger $n$),
    and available compute budget.  Adaptive selection analogous to scree plots in PCA
    or BIC in model selection does not exist for spatial subsampling.

  \item[\textbf{Approach}]
    Run REX-SUB for a sequence $n = n_{\min}, n_{\min}+\Delta, \ldots, n_{\max}$
    using previous solutions as warm starts (reducing marginal cost).  At each $n$,
    record MSPE on a held-out test set.  Fit a parametric decay model
    $\widehat{\text{MSPE}}(n) = a\,n^{-\beta} + c$ and identify the elbow via the
    point of maximum curvature.  Alternatively, stop when
    $|\Delta\widehat{\text{MSPE}}| < \varepsilon$ for a user-specified tolerance.

  \item[\textbf{Key experiment}]
    Across 8 simulation settings (Table~1 of the paper), generate MSPE($n$) curves
    for $n \in \{10, 15, 25, 36, 50, 75, 100, 150, 200\}$.  Characterise how the
    elbow location depends on $\nu$, $\rho^*$, and noise.

  \item[\textbf{Success criterion}]
    Adaptive $n$-selection matches the MSPE-optimal $n$ (chosen by oracle knowledge
    of the test set) within $\pm 10$ on $\geq 12$ of 16 simulation settings.

  \item[\novelty{4}]
    Adaptive subsample size for GP subsampling is entirely unaddressed.

  \item[\feasibility{5}]
    Computationally straightforward; warm-starting makes the sequence affordable.
\end{description}

\end{ideabox}

% ── Idea 3 ───────────────────────────────────────────────────────────────────
\subsection{Idea 3: REX-SUB-ST --- Spatio-Temporal Subsampling}

\begin{ideabox}{Idea~3 \textbullet{} Model Extension [Gap~G2]}

\begin{description}[leftmargin=5em, labelwidth=4.5em]
  \item[\textbf{Hypothesis}]
    Extending REX-SUB to spatio-temporal GPs by treating $(\mathbf{s}_i, t_j)$
    pairs as candidate subsample points will yield lower MSPE on multi-temporal
    satellite data than applying REX-SUB independently to each time snapshot.

  \item[\textbf{Motivation}]
    The MODIS satellite collects precipitable water data daily; the paper uses only
    June 22, 2023.  Atmospheric and environmental processes exhibit strong temporal
    autocorrelation --- subsampling across time is not equivalent to independently
    subsampling $T$ spatial snapshots.  No spatio-temporal optimal subsampling
    method exists.  Nonseparable covariance functions (Gneiting, 2002) and
    spatio-temporal \Vecchia{} approximations (Guinness, 2018) provide the necessary
    modelling infrastructure.

  \item[\textbf{Approach}]
    (i) Index candidate points as $(\mathbf{s}_i, t_j)$ pairs.
    (ii)~Use a spatio-temporal \Vecchia{} approximation with maxmin ordering adapted
    for space--time (ordering by $\|\mathbf{s}\| + \lambda t$ for a scale parameter
    $\lambda$).  (iii)~Define a spatio-temporal \mspe{} criterion over a held-out
    space--time test set.  (iv)~Run Algorithm~1 with swaps on $(\mathbf{s}, t)$ pairs.

  \item[\textbf{Key experiment}]
    Apply to 30 days of MODIS data ($N_{\text{total}} \approx 80$M).  Compare:
    (a) REX-SUB applied independently per day ($T \times n$ total subsample),
    (b) REX-SUB-ST jointly across days (same total subsample budget).

  \item[\textbf{Success criterion}]
    REX-SUB-ST achieves lower spatio-temporal MSPE than $T$ independent REX-SUB
    runs at the same total subsample size for $\geq 80\%$ of held-out space--time
    locations.

  \item[\novelty{5}]
    First optimal subsampling method for spatio-temporal GP models.

  \item[\feasibility{3}]
    Spatio-temporal \Vecchia{} exists; the main challenge is the larger candidate
    space and higher GP fitting cost per exchange.
\end{description}

\end{ideabox}

% ── Idea 4 ───────────────────────────────────────────────────────────────────
\subsection{Idea 4: REX-SUB-NS --- Non-Stationary Spatial Subsampling}

\begin{ideabox}{Idea~4 \textbullet{} Model Extension [Gap~G1]}

\begin{description}[leftmargin=5em, labelwidth=4.5em]
  \item[\textbf{Hypothesis}]
    Replacing the global stationary GP criterion in REX-SUB with a locally adaptive
    non-stationary criterion will yield lower MSPE on spatially heterogeneous fields
    while selecting subsample locations that cluster in high-variation regions.

  \item[\textbf{Motivation}]
    REX-SUB fits a single Matérn covariance to the entire domain.  On non-stationary
    fields, this misspecification causes the exchange algorithm to either under-sample
    high-variation regions (if global $\phi$ is over-estimated) or over-sample smooth
    regions (if $\phi$ is under-estimated).  Locally stationary GP approximations
    such as the scalable partitioned model of Lee \& Park (2023) can characterise
    local spatial structure, providing a heterogeneous criterion for the exchange.

  \item[\textbf{Approach}]
    Replace the Vecchia-fitted global GP within Algorithm~1 with a partition-based
    locally stationary GP (e.g.\ Lee \& Park, 2023).  Define the \mspe{} criterion
    as a spatially weighted average $\phi_{\text{NS}} = \sum_i w(\mathbf{s}_i)
    (\hat{Z}(\mathbf{s}_i) - Z(\mathbf{s}_i))^2$ where $w(\mathbf{s})$ reflects
    local variance.  Alternatively, fit separate covariance parameters per partition
    and propagate these through the exchange.

  \item[\textbf{Key experiment}]
    Simulate a non-stationary field with known spatially varying range parameter
    $\phi(\mathbf{s})$ (e.g.\ using the SPDE approach or convoSPAT).  Compare
    (a) REX-SUB (stationary criterion), (b) REX-SUB-NS (locally adaptive criterion),
    and (c) LHS on MSPE separately in smooth vs.\ rough regions.

  \item[\textbf{Success criterion}]
    REX-SUB-NS achieves $\geq 15\%$ lower MSPE in non-stationary regions compared
    to standard REX-SUB, with computation time $\leq 2\times$ the standard method.

  \item[\novelty{5}]
    No subsampling method accounts for spatial non-stationarity in its criterion.

  \item[\feasibility{3}]
    Non-stationary GP fitting within the exchange loop increases complexity; local
    parameter estimation must remain scalable.
\end{description}

\end{ideabox}

% ── Idea 5 ───────────────────────────────────────────────────────────────────
\subsection{Idea 5: REX-SUB-NG --- Non-Gaussian Spatial Responses}

\begin{ideabox}{Idea~5 \textbullet{} Model Extension [Gap~G6]}

\begin{description}[leftmargin=5em, labelwidth=4.5em]
  \item[\textbf{Hypothesis}]
    Adapting REX-SUB's optimality criterion to the Continuous Ranked Probability
    Score (CRPS) under a latent GP with non-Gaussian likelihood will enable
    prediction-optimal subsampling for spatial count, binary, and ordinal data.

  \item[\textbf{Motivation}]
    Many environmental and ecological spatial datasets are non-Gaussian: bird species
    counts (Poisson), disease incidence rates (negative binomial), habitat
    presence/absence (Bernoulli), and soil quality categories (ordinal).
    Log-transforming count data to approximate Gaussianity, as done in the paper,
    is valid only for certain distributions and may distort prediction intervals.
    Zilber \& Katzfuss (2021) provide a Vecchia-Laplace approximation for
    non-Gaussian spatial likelihoods --- this can serve as a direct replacement for
    the Gaussian Vecchia within Algorithm~1.

  \item[\textbf{Approach}]
    Replace the Vecchia-Gaussian likelihood with Vecchia-Laplace (Zilber \&
    Katzfuss, 2021) in Algorithm~1.  Replace $\phi_{\text{MSPE}}$ with
    $\phi_{\text{CRPS}} = \frac{1}{|I_{\text{test}}|}\sum_{i} \text{CRPS}(F_i,
    Z(\mathbf{s}_i))$ where $F_i$ is the predictive CDF at test location
    $\mathbf{s}_i$.  Maintain the same exchange structure.

  \item[\textbf{Key experiment}]
    Apply to the North American Breeding Bird Survey dataset (species counts across
    spatial locations in North America).  Compare REX-SUB-NG (CRPS, Vecchia-Laplace)
    vs.\ REX-SUB on log-counts vs.\ LHS on held-out log-score.

  \item[\textbf{Success criterion}]
    REX-SUB-NG achieves lower predictive CRPS than all baselines on $\geq 3$
    species with markedly non-Gaussian count distributions.

  \item[\novelty{5}]
    No prediction-optimal subsampling exists for non-Gaussian spatial GP models.

  \item[\feasibility{3}]
    Vecchia-Laplace integration is more complex; gradient computation for the
    Laplace approximation within the exchange adds overhead.
\end{description}

\end{ideabox}

% ── Idea 6 ───────────────────────────────────────────────────────────────────
\subsection{Idea 6: Theoretical Analysis --- Convergence and Consistency}

\begin{ideabox}{Idea~6 \textbullet{} Theory [Gap~G8]}

\begin{description}[leftmargin=5em, labelwidth=4.5em]
  \item[\textbf{Hypothesis}]
    Under regularity conditions on the spatial covariance function and spatial
    domain filling, the predictions from a REX-SUB-selected subsample converge
    to the full-data GP predictions at a characterisable rate in $n$, $N$, $\phi$,
    and $\nu$.

  \item[\textbf{Motivation}]
    REX-SUB is entirely empirical.  IBOSS (Wang et al., 2019) shows that for linear
    regression, D-optimal subsampling achieves asymptotic efficiency equal to the
    full dataset.  Three theoretical questions are open for REX-SUB:
    (i)~Is the exchange-algorithm solution consistent for the oracle
    (IMSPE-minimising) subsample as $n_{\text{cand}}, n_{\text{repeat}} \to \infty$?
    (ii)~What is the convergence rate of subsample-GP MSPE to full-data GP MSPE as
    $n \to N$?
    (iii)~Are covariance parameter estimates from the informative subsample
    asymptotically unbiased?

  \item[\textbf{Approach}]
    Frame REX-SUB as a stochastic approximation to the MSPE-minimising subsampling
    problem.  Use spatial mixing conditions (strong mixing or $m$-dependence) to
    handle correlated observations.  Derive an upper bound on
    $\mathbb{E}[\text{MSPE}_{\text{subsample}}] - \text{MSPE}_{\text{oracle}}$
    as a function of $n$, $N$, $\phi$, $\nu$, and the exchange stopping criteria.
    For parameter bias, use the theory of informative sampling (cf.\ Pfeffermann,
    1993) adapted to the spatial GP setting.

  \item[\textbf{Key experiment}]
    Empirically validate derived rates by running REX-SUB across
    $n \in \{10, 25, 50, 100, 250\}$ for large $N = 100{,}000$ simulated
    datasets.  Plot observed MSPE decay rate against theoretical predictions.

  \item[\textbf{Success criterion}]
    Derived upper bound matches observed convergence rate within a constant factor
    across $\geq 6$ of 8 simulation settings.

  \item[\novelty{5}]
    No convergence theory for prediction-optimal spatial subsampling exists.

  \item[\feasibility{2}]
    Requires significant probabilistic spatial analysis; hardest idea to execute
    but highest long-term scientific impact.
\end{description}

\end{ideabox}

% ─────────────────────────────────────────────────────────────────────────────
\section{Prioritised Research Roadmap}
% ─────────────────────────────────────────────────────────────────────────────

\begin{table}[ht]
\centering
\caption{Prioritised research extension roadmap.  Effort estimates are for a
  single PhD student or postdoc.  Dependencies are noted where one idea unlocks
  another.}
\label{tab:roadmap}
\small
\setlength{\tabcolsep}{4pt}
\begin{tabular}{@{}clccp{3.8cm}c@{}}
\toprule
\textbf{Priority} & \textbf{Idea} & \textbf{Novelty} & \textbf{Feasibility}
  & \textbf{Justification} & \textbf{Effort} \\
\midrule
\rowcolor{ideabg}
1 & Idea~3: Rank-1 Vecchia Updates & 4/5 & 4/5
  & Unlocks all other extensions by making large $n$ practical;
    algebraically well-scoped
  & 3--6 mo. \\
\addlinespace[2pt]
2 & Idea~2: Adaptive subsample size $n$ & 4/5 & 5/5
  & Resolves most common practitioner question; computationally
    straightforward with warm-starting
  & 2--4 mo. \\
\addlinespace[2pt]
\rowcolor{ideabg}
3 & Idea~1: REX-SUB-NS (non-stationary) & 5/5 & 3/5
  & Most restrictive assumption; highest practical relevance for
    real environmental data
  & 6--12 mo. \\
\addlinespace[2pt]
4 & Idea~3: REX-SUB-ST (spatio-temporal) & 5/5 & 3/5
  & Natural next application given the MODIS motivation;
    \textit{requires Priority~1}
  & 6--12 mo. \\
\addlinespace[2pt]
\rowcolor{ideabg}
5 & Idea~5: REX-SUB-NG (non-Gaussian) & 5/5 & 3/5
  & Opens entirely new application domains (ecology, epidemiology);
    \textit{requires Priority~1}
  & 6--12 mo. \\
\addlinespace[2pt]
6 & Idea~6: Convergence theory & 5/5 & 2/5
  & Highest scientific prestige; needed for the methodology to be
    a lasting contribution; \textit{benefits from Priorities 1--5}
  & 12--24 mo. \\
\bottomrule
\end{tabular}
\end{table}

% ─────────────────────────────────────────────────────────────────────────────
\section{Cross-Cutting Insight}
% ─────────────────────────────────────────────────────────────────────────────

\begin{critbox}[title=\textbf{The Single Highest-Leverage Gap}]

The deepest structural gap --- cutting across all six extension directions --- is
the \textbf{absence of a fast criterion update rule} (Gap~G3, Idea~1).

Every proposed extension (non-stationary, spatio-temporal, non-Gaussian, larger $n$)
is computationally gated by the cost of re-fitting the GP at each exchange step.
Table~5 of the paper shows that for $n_{\text{cand}} = 10$ and
$n_{\text{repeat}} = 2$, REX-SUB fits
$n_{\text{repeat}} \times n \times n_{\text{cand}} = 2 \times 25 \times 10 = 500$
GPs per subsample.  Rank-1 Vecchia updates would reduce each of these 500 fits to a
$\mathcal{O}(m^2)$ matrix update rather than a full $\mathcal{O}(m^2 n)$
re-factorisation --- a factor of $n$ speedup.

\textbf{Idea~1 is therefore the recommended first contribution}: it is a prerequisite
that multiplicatively accelerates all subsequent extensions and transforms REX-SUB
from a method feasible for $n \leq 36$ into one usable for $n \leq 500$.

\end{critbox}

\medskip

\begin{warnbox}
  \textbf{Second structural gap:}
  The stationarity assumption (Gap~G1) makes REX-SUB unreliable for the most common
  real-world spatial datasets.  The strong spatial heterogeneity visible in Figure~2
  of the paper (smooth Atlantic Ocean vs.\ complex land--ocean boundary with high
  precipitable water contrast) suggests the method's MODIS success may be partially
  coincidental --- the log-transformation and chosen geographic extent may happen to
  produce approximately stationary residuals.  REX-SUB-NS (Idea~4) addresses this
  directly.
\end{warnbox}

% ─────────────────────────────────────────────────────────────────────────────
\section*{Acknowledgements}
% ─────────────────────────────────────────────────────────────────────────────

This gap analysis was produced using the AI Auto-Research Skill Suite, implementing
the research lifecycle framework described in:

\begin{quote}
  Kong, L., Sun, X., Chow, W., et al.\ (2026).
  \emph{AI for Auto-Research: Roadmap \& User Guide.}
  arXiv:2605.18661v1 [cs.AI], 18 May 2026.
\end{quote}

Skills applied: \texttt{/literature-review} (Phase~1, Stage~2),
\texttt{/peer-review} (Phase~3, Stage~6), and
\texttt{/idea-generation} (Phase~1, Stage~1).

% ─────────────────────────────────────────────────────────────────────────────
\section*{References}
% ─────────────────────────────────────────────────────────────────────────────

\begin{thebibliography}{99}

\bibitem{Banerjee2008}
Banerjee, S., Gelfand, A.E., Finley, A.O., Sang, H. (2008).
Gaussian predictive process models for large spatial data sets.
\textit{Journal of the Royal Statistical Society: Series B}, 70(4), 825--848.

\bibitem{Binois2021}
Binois, M., Gramacy, R.B. (2021).
hetGP: Heteroskedastic Gaussian process modeling and sequential design in R.
\textit{Journal of Statistical Software}, 98(13), 1--44.

\bibitem{Cheng2020}
Cheng, Q., Wang, H., Yang, M. (2020).
Information-based optimal subdata selection for big data logistic regression.
\textit{Journal of Statistical Planning and Inference}, 209, 112--122.

\bibitem{Cressie1999}
Cressie, N., Huang, H.C. (1999).
Classes of nonseparable, spatio-temporal stationary covariance functions.
\textit{Journal of the American Statistical Association}, 94(448), 1330--1339.

\bibitem{Cressie2008}
Cressie, N., Johannesson, G. (2008).
Fixed rank kriging for very large spatial data sets.
\textit{Journal of the Royal Statistical Society: Series B}, 70(1), 209--226.

\bibitem{CressieWikle2011}
Cressie, N., Wikle, C.K. (2011).
\textit{Statistics for Spatio-Temporal Data}.
John Wiley \& Sons.

\bibitem{Datta2016}
Datta, A., Banerjee, S., Finley, A.O., Gelfand, A.E. (2016).
Hierarchical nearest-neighbor Gaussian process models for large geostatistical datasets.
\textit{Journal of the American Statistical Association}, 111(514), 800--812.

\bibitem{Fuentes2001}
Fuentes, M. (2001).
A high frequency kriging approach for non-stationary environmental processes.
\textit{Environmetrics}, 12(5), 469--483.

\bibitem{Furrer2006}
Furrer, R., Genton, M.G., Nychka, D. (2006).
Covariance tapering for interpolation of large spatial datasets.
\textit{Journal of Computational and Graphical Statistics}, 15(3), 502--523.

\bibitem{Gneiting2002}
Gneiting, T. (2002).
Nonseparable, stationary covariance functions for space--time data.
\textit{Journal of the American Statistical Association}, 97(458), 590--600.

\bibitem{Guinness2018}
Guinness, J. (2018).
Permutation and grouping methods for sharpening Gaussian process approximations.
\textit{Technometrics}, 60(4), 415--429.

\bibitem{Hayashi2020}
Hayashi, K., Imaizumi, M., Yoshida, Y. (2020).
On random subsampling of Gaussian process regression.
\textit{PMLR}, pp.\ 2055--2065.

\bibitem{Heaton2019}
Heaton, M.J., Datta, A., Finley, A.O., et al.\ (2019).
A case study competition among methods for analyzing large spatial data.
\textit{Journal of Agricultural, Biological and Environmental Statistics},
24(3), 398--425.

\bibitem{Hefley2017}
Hefley, T.J., Broms, K.M., Brost, B.M., et al.\ (2017).
The basis function approach for modeling autocorrelation in ecological data.
\textit{Ecology}, 98(3), 632--646.

\bibitem{Katzfuss2017}
Katzfuss, M. (2017).
A multi-resolution approximation for massive spatial datasets.
\textit{Journal of the American Statistical Association}, 112(517), 201--214.

\bibitem{Katzfuss2020}
Katzfuss, M., Guinness, J., Gong, W., Zilber, D. (2020).
Vecchia approximations of Gaussian-process predictions.
\textit{Journal of Agricultural, Biological and Environmental Statistics},
25(3), 383--414.

\bibitem{Kim2005}
Kim, H.M., Mallick, B.K., Holmes, C.C. (2005).
Analyzing nonstationary spatial data using piecewise Gaussian processes.
\textit{Journal of the American Statistical Association}, 100(470), 653--668.

\bibitem{Lee2023}
Lee, B.S., Park, J. (2023).
A scalable partitioned approach to model massive nonstationary non-Gaussian spatial datasets.
\textit{Technometrics}, 65(1), 105--116.

\bibitem{Lindgren2011}
Lindgren, F., Rue, H., Lindström, J. (2011).
An explicit link between Gaussian fields and Gaussian Markov random fields.
\textit{Journal of the Royal Statistical Society: Series B}, 73(4), 423--498.

\bibitem{Liu2023}
Liu, Y., Stufken, J., Yang, M. (2023).
Information-based optimal subdata selection for clusterwise linear regression.
arXiv preprint arXiv:2309.00720.

\bibitem{McKay1979}
McKay, M.D., Beckman, R.J., Conover, W.J. (1979).
A comparison of three methods for selecting values of input variables.
\textit{Technometrics}, 21(2), 239--245.

\bibitem{Rasmussen2005}
Rasmussen, C.E., Williams, C.K.I. (2005).
\textit{Gaussian Processes for Machine Learning}.
The MIT Press.

\bibitem{Risser2015}
Risser, M.D., Calder, C.A. (2015).
Local likelihood estimation for covariance functions with spatially varying parameters.
arXiv preprint arXiv:1507.08613.

\bibitem{Saha2024}
Saha, A., Bradley, J.R. (2024).
Incorporating subsampling into Bayesian models for high-dimensional spatial data.
\textit{Bayesian Analysis}, 1(1), 1--40.

\bibitem{Sang2011}
Sang, H., Jun, M., Huang, J.Z. (2011).
Covariance approximation for large multivariate spatial data sets with an application
to multiple climate model errors.
\textit{Annals of Applied Statistics}, 5(4), 2494--2526.

\bibitem{ShiTang2021}
Shi, C., Tang, B. (2021).
Model-robust subdata selection for big data.
\textit{Journal of Statistical Theory and Practice}, 15, 1--17.

\bibitem{Singh2023}
Singh, B., Stufken, J. (2023).
Subdata selection with a large number of variables.
\textit{New England Journal of Statistical Data Science}, 1(3), 426--438.

\bibitem{Vecchia1988}
Vecchia, A.V. (1988).
Estimation and model identification for continuous spatial processes.
\textit{Journal of the Royal Statistical Society: Series B}, 50, 297--312.

\bibitem{Wang2019}
Wang, H., Stufken, J., Yang, M. (2019).
Information-based optimal subdata selection for big data linear regression.
\textit{Journal of the American Statistical Association}, 114(525), 393--405.

\bibitem{Zilber2021}
Zilber, D., Katzfuss, M. (2021).
Vecchia--Laplace approximations of generalised Gaussian processes for big
non-Gaussian spatial data.
\textit{Computational Statistics and Data Analysis}, 153, 107081.

\bibitem{ZammitMangion2021}
Zammit-Mangion, A., Cressie, N. (2021).
FRK: An R package for spatial and spatio-temporal prediction with large datasets.
\textit{Journal of Statistical Software}, 98(4), 1--48.

\end{thebibliography}

\end{document}
